\section{第11章 ソフトウェア工学モデルと方法}
\noindent この資料の主張\par
ソフトウェア工学モデルと方法は、開発作業を場当たり的な作業から、説明可能で、反復可能で、改善しやすい作業へ変えるための道具である。モデルは、対象をすべて写し取るものではなく、意思決定に必要な側面を抽象化した表現である。方法は、要求、設計、実装、テスト、検証を体系的に進めるための手順である。したがって、この章の中心は「何をモデル化し、どう読み、どう確かめ、どの方法で進めるか」を判断できるようになることである。

\begin{pointbox}{この資料の読み方}
専門用語は原則として日本語で示し、必要に応じて英語を括弧内に併記する。英語名は、元資料、ツール、論文、標準で確認するときの手がかりとして残す。初学者が前提知識なしに読めるように、各節では「何のための概念か」「どの失敗を防ぐのか」「実務では何を見るか」を重視する。文体は、だである調で統一する。
\end{pointbox}
